\section{深度 Q 网络变体与训练对比}

\subsection{关键变体剖析}

\begin{table}[H]
\centering
\begin{tabular}{lll}
\toprule
变体 & 核心机制 & 影响 \\
\midrule
Double DQN & 当前网络选动作，目标网络评估 & 抑制 overestimation，稳定性提升 \\
Dueling DQN & 分解 $Q=V+A$ & 更快学到状态价值，动作优势可共享参数 \\
Prioritized Replay & 根据 TD error 加权采样 & 快速拾取稀有 but 关键的转移 \\
Rainbow & 多个增强技巧组合（NoisyNet、Multi-step、Distributional） & 覆盖性更广，但调参更复杂 \\
\bottomrule
\end{tabular}
\caption{DQN 变体与其工程影响}
\end{table}

\begin{knowledgebox}{多维组合的适用场景}
不是所有场景都需要 Rainbow 的全套增强。在低容量环境下，Double DQN 结合 Dueling 已足以抑制过估。在富有计算资源的 Atari 任务或模拟驾驶中，Rainbow 多步回报与 NoisyNet 提供更强的泛化能力。
\end{knowledgebox}

\subsection{训练与推理成本对比}

\begin{importantbox}{训练 vs 推理 成本模型}
训练集中于采样与优化：$Cost_{train} \approx N_{\text{steps}} \times (^{\text{forward}}C + ^{\text{backward}}C)$；推理则在吞吐与延迟间权衡：$Cost_{infer} \approx C_{\text{forward}} + C_{\text{action selection}}$。若将 DQN 推向嵌入式设备，应考虑量化／剪枝以减少计算量。
\end{importantbox}

\begin{knowledgebox}{推理栈拆解}
Inference pipeline 通常包括：环境感知（state encoding）、Q 网络前向、$\arg\max$ 动作选择、行动滤波（如先验约束）、以及与模拟器接口的安全检查。在多 agent 协同中，还需把握通信与同步延迟。
\end{knowledgebox}

\begin{warningbox}{低延迟推理中的过度量化}
为了优化推理 latency，工程团队往往会对网络进行量化或剪枝，但如果没有保留关键路径（如最后一层 advantage 分支），Q 值排序可能被破坏。推荐保持双通道（量化用于推理，full precision 备份用于 periodic calibration）。
\end{warningbox}

\subsection{本章小结}

本章通过 DQN 变体对比与成本模型讨论，指出在不同资源约束下如何选择合适的增强技巧与算子组合。

